\documentclass[12pt]{article}

\usepackage[a4paper, top=2cm, bottom=2cm, left=2.5cm, right=2.5cm]{geometry}
\usepackage{amsmath, amssymb}
\usepackage{fontspec}
\usepackage{mathastext}
\usepackage{enumitem}
\usepackage{array}
\usepackage{tabularx}
\usepackage[table]{xcolor}
\usepackage{booktabs}
\usepackage{hyperref}
\usepackage{bookmark}

% ---------------- FONTS ----------------
\setmainfont[
    UprightFont = Handlee-Regular.ttf,
    AutoFakeBold = 2.5
]{Handlee}
\newfontfamily\headingfont{PT.ttf}

% ---------------- HANDWRITTEN MATH ----------------
\Mathastext
\everymath{\displaystyle}

% ---------------- GENERAL ----------------
\setlength{\parindent}{0pt}
\pagenumbering{gobble}
\linespread{1.2}

\newcommand{\mychapter}[2]{%
    \phantomsection
    \hypertarget{#2}{}%
    \bookmark[level=0,dest=#2]{#1}%
    \begin{center}
        {\headingfont\Huge #1}
        \par
        \vspace{4pt}
        \hrule
    \end{center}
}

\newcommand{\mysection}[3]{%
    \phantomsection
    \hypertarget{#3}{}%
    \bookmark[level=1,dest=#3]{#1\space #2}%
    \newpage
    \noindent{\headingfont\Large #1\space #2\par}
    \vspace{4pt}
    \hrule
    \vspace{10pt}
}

\newcommand{\mysubsection}[2]{%
    \noindent{\headingfont\large #1\space #2\par}
    \vspace{4pt}
    \hrule
    \vspace{8pt}
}

\begin{document}

\mychapter{Chapter 1 : Introduction \& Basic Concepts}{Chapter 1 : Introduction \& Basic Concepts}

\begin{center}
\begin{tabular}{|l|l|}
\hline
\textbf{Section} & \textbf{Core Revision Focus} \\
\hline
1.1 & Thermodynamics, energy, first/second laws, approaches \\
\hline
1.2 & Dimensions, units, dimensional homogeneity, conversions \\
\hline
1.3 & Systems, control volumes, mass/energy transfer \\
\hline
1.4 & Properties, intensive/extensive, specific properties, continuum \\
\hline
1.5 & Density, specific volume, SG, specific weight \\
\hline
1.6 & State, equilibrium, state postulate \\
\hline
1.7 & Processes, paths, cycles, steady flow \\
\hline
1.8 & Temperature, zeroth law, temperature scales \\
\hline
1.9 & Pressure, pressure relations, hydrostatics, Pascal's law \\
\hline
1.10 & Barometer, manometer, differential pressure \\
\hline
1.11 & Systematic engineering problem solving \\
\hline
\end{tabular}
\end{center}

\newpage

%=========================================================
\mysection{1.1}{Thermodynamics and Energy}{1.1}

\mysubsection{}{Core Concepts}

\begin{itemize}
    \item \textbf{Thermodynamics}: science of energy, energy transformations, and interactions between energy and matter.
    \item Energy may be viewed as the \textbf{ability to cause changes}.
    \item \textbf{Conservation of energy}: energy cannot be created or destroyed; only transformed.
\end{itemize}

\begin{equation*}
    \boxed{\Delta E=E_{\mathrm{in}}-E_{\mathrm{out}}}
\end{equation*}

\begin{center}
\begin{tabular}{|c|c|}
\hline
\textbf{Law} & \textbf{Main principle} \\
\hline
First law & Energy is conserved; energy is a thermodynamic property. \\
\hline
Second law & Energy has both quantity and quality; actual processes proceed toward lower energy quality. \\
\hline
\end{tabular}
\end{center}

\mysubsection{}{Thermodynamic Approaches}

\begin{center}
\begin{tabular}{|c|c|c|}
\hline
\textbf{Approach} & \textbf{Basis} & \textbf{Feature} \\
\hline
Classical & Macroscopic behavior & Direct, relatively simple \\
\hline
Statistical & Average microscopic behavior & More involved \\
\hline
\end{tabular}
\end{center}

\begin{itemize}
    \item \textbf{Classical thermodynamics}: uses measurable macroscopic quantities; primary engineering approach.
    \item \textbf{Statistical thermodynamics}: relates macroscopic behavior to microscopic particle behavior.
\end{itemize}

\textbf{Key second-law idea:} heat naturally transfers from higher temperature to lower temperature, with degradation of energy quality.

%=========================================================
\mysection{1.2}{Importance of Dimensions and Units}{1.2}

\mysubsection{}{Dimensions and Units}

\begin{itemize}
    \item \textbf{Dimensions}: characteristics used to describe physical quantities.
    \item \textbf{Units}: magnitudes assigned to dimensions.
    \item \textbf{Primary dimensions}: fundamental quantities such as mass, length, time, temperature.
    \item \textbf{Secondary dimensions}: derived from primary dimensions, e.g. velocity, energy, volume.
\end{itemize}

\begin{center}
\begin{tabular}{|c|c|}
\hline
\textbf{System} & \textbf{Description} \\
\hline
SI & International System of Units; decimal-based \\
\hline
English/USCS & United States Customary System \\
\hline
\end{tabular}
\end{center}

\mysubsection{}{SI Fundamental Units}

\begin{center}
\begin{tabular}{|c|c|c|}
\hline
\textbf{Dimension} & \textbf{Unit} & \textbf{Symbol} \\
\hline
Length & meter & m \\
\hline
Mass & kilogram & kg \\
\hline
Time & second & s \\
\hline
Temperature & kelvin & K \\
\hline
Electric current & ampere & A \\
\hline
Amount of light & candela & cd \\
\hline
Amount of matter & mole & mol \\
\hline
\end{tabular}
\end{center}

\mysubsection{}{SI Prefixes}

\begin{center}
\begin{tabular}{|c|c|c@{\qquad}c|c|c|}
\hline
\textbf{Factor} & \textbf{Prefix} & \textbf{Symbol}
& \textbf{Factor} & \textbf{Prefix} & \textbf{Symbol} \\
\hline
$10^{24}$ & yotta & Y & $10^{-1}$ & deci & d \\
$10^{21}$ & zetta & Z & $10^{-2}$ & centi & c \\
$10^{18}$ & exa & E & $10^{-3}$ & milli & m \\
$10^{15}$ & peta & P & $10^{-6}$ & micro & $\mu$ \\
$10^{12}$ & tera & T & $10^{-9}$ & nano & n \\
$10^9$ & giga & G & $10^{-12}$ & pico & p \\
$10^6$ & mega & M & $10^{-15}$ & femto & f \\
$10^3$ & kilo & k & $10^{-18}$ & atto & a \\
$10^2$ & hecto & h & $10^{-21}$ & zepto & z \\
$10^1$ & deka & da & $10^{-24}$ & yocto & y \\
\hline
\end{tabular}
\end{center}

\mysubsection{}{Unit Conventions}

\begin{itemize}
    \item Unit names are lowercase: newton, pascal, joule.
    \item Symbols derived from proper names are capitalized: $N$, $Pa$, $J$.
    \item Unit symbols are not pluralized and normally have no period.
\end{itemize}

\mysubsection{}{Dimensional Homogeneity}

Every term in a valid engineering equation must have compatible dimensions/units.

\begin{itemize}
    \item Different units cannot be added or subtracted.
    \item Dimensional analysis can detect equation/calculation errors.
    \item Dimensional homogeneity is \textbf{necessary but not sufficient} for physical correctness.
\end{itemize}

\mysubsection{}{Unity Conversion Ratios}

A ratio of equivalent quantities is exactly one:

\begin{equation*}
\frac{1~\mathrm{N}}
{1~\mathrm{kg\,m/s^2}}=1
\end{equation*}

\begin{equation*}
\frac{1~\mathrm{lbf}}
{32.174~\mathrm{lbm\,ft/s^2}}=1
\end{equation*}

\begin{itemize}
    \item Insert unity ratios without changing physical value.
    \item Arrange ratios so unwanted units cancel.
    \item The inverse ratio is also equal to $1$.
    \item Preferred for unit conversion rather than introducing $g_c$ merely to force units to match.
\end{itemize}

%=========================================================
\mysection{1.3}{Systems and Control Volumes}{1.3}

\mysubsection{}{Basic Definitions}

\begin{center}
\begin{tabular}{|c|c|}
\hline
\textbf{Term} & \textbf{Definition} \\
\hline
System & Quantity of matter or region in space selected for study \\
\hline
Surroundings & Everything outside the system \\
\hline
Boundary & Real/imaginary surface separating system and surroundings \\
\hline
\end{tabular}
\end{center}

Boundary may be \textbf{fixed/movable} and \textbf{real/imaginary}; mathematically it has zero thickness.

\mysubsection{}{System Classification}

\begin{center}
\begin{tabular}{|c|c|c|}
\hline
\textbf{System} & \textbf{Mass Transfer} & \textbf{Energy Transfer} \\
\hline
Closed system & No & Yes \\
\hline
Isolated system & No & No \\
\hline
Open system & Yes & Yes \\
\hline
\end{tabular}
\end{center}

\begin{itemize}
    \item \textbf{Closed system/control mass}: fixed amount of mass; heat/work may cross boundary.
    \item \textbf{Isolated system}: special closed system with neither mass nor energy crossing.
    \item \textbf{Open system/control volume}: selected region in space through which mass and energy may cross.
\end{itemize}

For a closed system:
\begin{equation*}
m=\text{constant}
\end{equation*}

\mysubsection{}{Control Volume}

\begin{itemize}
    \item Boundary of a control volume = \textbf{control surface}.
    \item Control surface may be real/imaginary and fixed/moving.
    \item Proper control-volume selection can greatly simplify analysis.
    \item Typical devices: compressors, turbines, nozzles, pumps, boilers.
\end{itemize}

\begin{center}
\begin{tabular}{|c|c|c|}
\hline
\textbf{Feature} & \textbf{Closed System} & \textbf{Control Volume} \\
\hline
Matter & Fixed amount & May change \\
\hline
Mass crossing boundary & No & Yes \\
\hline
Energy crossing boundary & Yes & Yes \\
\hline
Boundary & Real/imaginary & Real/imaginary \\
\hline
Boundary motion & Fixed/movable & Fixed/movable \\
\hline
Typical example & Piston-cylinder & Turbine/nozzle \\
\hline
\end{tabular}
\end{center}

%=========================================================
\mysection{1.4}{Properties of a System}{1.4}

\mysubsection{}{Intensive and Extensive Properties}

\begin{itemize}
    \item \textbf{Property}: any characteristic of a system, e.g. $P,T,V,m$.
    \item \textbf{Intensive}: independent of system size/mass.
    \item \textbf{Extensive}: depends on system size/mass.
\end{itemize}

\begin{center}
\begin{tabular}{|c|c|}
\hline
\textbf{Type} & \textbf{Examples} \\
\hline
Intensive & $T$, $P$, density \\
\hline
Extensive & $m$, $V$, total momentum \\
\hline
\end{tabular}
\end{center}

\textbf{Division test:}
\begin{itemize}
    \item Intensive $\rightarrow$ value remains unchanged.
    \item Extensive $\rightarrow$ value changes in proportion to amount of matter.
\end{itemize}

\mysubsection{}{Specific Properties}

An extensive property per unit mass is a \textbf{specific property}; hence it is intensive.

\begin{equation*}
\boxed{v=\frac{V}{m}}
\qquad
\boxed{e=\frac{E}{m}}
\end{equation*}

where $v$ = specific volume and $e$ = specific total energy.

\mysubsection{}{Continuum}

\begin{itemize}
    \item Matter is treated as a continuous, homogeneous substance rather than discrete molecules.
    \item Allows properties to vary continuously and be treated as point functions.
    \item Valid when system characteristic dimension $\gg$ molecular mean free path.
\end{itemize}

For oxygen at $1$ atm and $20^\circ\mathrm{C}$:
\begin{equation*}
d_{\mathrm{O_2}}\approx3\times10^{-10}\,\mathrm{m},
\qquad
\lambda\approx6.3\times10^{-8}\,\mathrm{m}
\end{equation*}

At very high vacuum/elevation, $\lambda$ may become comparable to system dimensions $\Rightarrow$ continuum assumption may fail.

%=========================================================
\mysection{1.5}{Density \& Specific Gravity}{1.5}

\mysubsection{}{Density and Specific Volume}

\begin{equation*}
\boxed{\rho=\frac{m}{V}}
\qquad
\boxed{v=\frac{V}{m}=\frac{1}{\rho}}
\end{equation*}

\begin{equation*}
\boxed{\rho v=1}
\end{equation*}

\begin{center}
\begin{tabular}{|c|c|c|}
\hline
\textbf{Quantity} & \textbf{Definition} & \textbf{SI Unit} \\
\hline
Density, $\rho$ & $m/V$ & $\mathrm{kg/m^3}$ \\
\hline
Specific volume, $v$ & $V/m$ & $\mathrm{m^3/kg}$ \\
\hline
\end{tabular}
\end{center}

For differential elements:
\begin{equation*}
\rho=\frac{\delta m}{\delta V}
\end{equation*}

\begin{itemize}
    \item Gas density generally $\uparrow$ with pressure and $\downarrow$ with temperature.
    \item Liquids/solids are essentially incompressible; density is much less pressure-sensitive.
\end{itemize}

\mysubsection{}{Specific Gravity}

\begin{equation*}
\boxed{SG=\frac{\rho}{\rho_{\mathrm{H_2O}}}}
\end{equation*}

For water at $4^\circ\mathrm{C}$:
\begin{equation*}
\rho_{\mathrm{H_2O}}=1000~\mathrm{kg/m^3}
\end{equation*}

$SG$ is \textbf{dimensionless}.

\begin{center}
\begin{tabular}{|c|c|}
\hline
$SG<1$ & Less dense than water \\
\hline
$SG=1$ & Same density as water \\
\hline
$SG>1$ & Denser than water \\
\hline
\end{tabular}
\end{center}

Useful relation:
\begin{equation*}
1~\mathrm{g/cm^3}
=1~\mathrm{kg/L}
=1000~\mathrm{kg/m^3}
\end{equation*}

Hence the numerical value of $SG$ equals density in $\mathrm{g/cm^3}$ or $\mathrm{kg/L}$.

Example:
\begin{equation*}
SG_{\mathrm{Hg}}=13.6
\Rightarrow
\rho_{\mathrm{Hg}}=13600~\mathrm{kg/m^3}
\end{equation*}

\mysubsection{}{Specific Weight}

\begin{equation*}
\boxed{\gamma_s=\rho g}
\end{equation*}

SI unit:
\begin{equation*}
[\gamma_s]=\mathrm{N/m^3}
\end{equation*}

%=========================================================
\mysection{1.6}{State and Equilibrium}{1.6}

\mysubsection{}{State and Equilibrium}

\begin{itemize}
    \item \textbf{State}: condition of a system described by its property values.
    \item Change in any property $\Rightarrow$ different state.
    \item \textbf{Equilibrium}: state of balance with no unbalanced potentials/driving forces.
    \item Thermodynamic equilibrium requires all relevant equilibrium conditions simultaneously.
\end{itemize}

\begin{center}
\begin{tabular}{|c|c|}
\hline
\textbf{Equilibrium} & \textbf{Condition} \\
\hline
Thermal & No temperature difference \\
\hline
Mechanical & No pressure changes with time at any point \\
\hline
Phase & Mass of each phase remains at equilibrium level \\
\hline
Chemical & Chemical composition does not change with time \\
\hline
\end{tabular}
\end{center}

\mysubsection{}{State Postulate}

For a \textbf{simple compressible system}:

\begin{equation*}
\boxed{\text{State is completely specified by two independent intensive properties.}}
\end{equation*}

Simple compressible system: electrical, magnetic, gravitational, motion, and surface-tension effects are absent/negligible.

\begin{itemize}
    \item $T$ and $v$ can specify the state.
    \item For single-phase systems, $T$ and $P$ are independent.
    \item For multiphase systems, $T$ and $P$ can be dependent:
\end{itemize}

\begin{equation*}
T=f(P)
\end{equation*}

Thus, $T$ and $P$ alone are insufficient to specify a two-phase state.

If another significant effect exists, an additional property is required for each important effect.

%=========================================================
\mysection{1.7}{Processes and Cycles}{1.7}

\mysubsection{}{Process}

\begin{itemize}
    \item \textbf{Process}: change from one equilibrium state to another.
    \item \textbf{Path}: series of equilibrium states through which the system passes.
    \item A process is characterized by initial state, final state, path, and system-surroundings interactions.
\end{itemize}

\mysubsection{}{Quasi-Equilibrium Process}

\begin{equation*}
\boxed{\text{System remains infinitesimally close to equilibrium throughout the process.}}
\end{equation*}

\begin{itemize}
    \item Approximated by a sufficiently slow process.
    \item Properties remain nearly uniform.
    \item Easier to analyze.
    \item Work-producing devices deliver maximum work under quasi-equilibrium operation.
\end{itemize}

For quasi-equilibrium processes, every point on the path represents an equilibrium state.

For nonquasi-equilibrium processes, a path for the system as a whole cannot be defined; it may be represented by a dashed line between initial and final states.

\mysubsection{}{Common Processes}

\begin{center}
\begin{tabular}{|c|c|}
\hline
\textbf{Process} & \textbf{Constant Property} \\
\hline
Isothermal & $T=\mathrm{constant}$ \\
\hline
Isobaric & $P=\mathrm{constant}$ \\
\hline
Isochoric/isometric & $v=\mathrm{constant}$ \\
\hline
\end{tabular}
\end{center}

\mysubsection{}{Cycle}

\begin{equation*}
\boxed{\text{Initial state}=\text{Final state}}
\end{equation*}

A cycle may contain several processes but returns the system to its initial state.

\mysubsection{}{Steady-Flow Process}

\begin{center}
\begin{tabular}{|c|c|}
\hline
Steady & No change with time \\
\hline
Uniform & No change with location \\
\hline
\end{tabular}
\end{center}

For a steady-flow control volume:
\begin{equation*}
\boxed{V=\mathrm{constant}},\qquad
\boxed{m=\mathrm{constant}},\qquad
\boxed{E=\mathrm{constant}}
\end{equation*}

At a fixed point, properties do not change with time, although they may vary from point to point.

Typical applications: turbines, pumps, boilers, and other continuously operating devices.

%=========================================================
\mysection{1.8}{Temperature and the Zeroth Law of Thermodynamics}{1.8}

\mysubsection{}{Thermal Equilibrium and Zeroth Law}

\begin{itemize}
    \item Heat transfers from higher temperature to lower temperature until temperatures become equal.
    \item Thermal equilibrium $\Rightarrow$ equal temperature.
\end{itemize}

\textbf{Zeroth law:}
\begin{equation*}
\boxed{
A\text{ in thermal equilibrium with }C,\;
B\text{ in thermal equilibrium with }C
\Rightarrow
A\text{ and }B\text{ are in thermal equilibrium.}
}
\end{equation*}

This provides the basis for temperature measurement using a thermometer as the third body.

\mysubsection{}{Temperature Scales}

\begin{center}
\begin{tabular}{|c|c|c|}
\hline
\textbf{Scale} & \textbf{Ice Point} & \textbf{Steam Point} \\
\hline
Celsius & $0^\circ\mathrm{C}$ & $100^\circ\mathrm{C}$ \\
\hline
Fahrenheit & $32^\circ\mathrm{F}$ & $212^\circ\mathrm{F}$ \\
\hline
\end{tabular}
\end{center}

\mysubsection{}{Absolute Temperature Scales}

\begin{itemize}
    \item SI thermodynamic scale: \textbf{Kelvin}, $K$.
    \item English thermodynamic scale: \textbf{Rankine}, $R$.
    \item Kelvin is written without degree symbol: $300~K$, not $300^\circ K$.
\end{itemize}

Absolute zero:
\begin{equation*}
\boxed{0~K=-273.15^\circ\mathrm{C}}
\end{equation*}

\mysubsection{}{Ideal-Gas Temperature Scale}

At sufficiently low pressure and constant volume:
\begin{equation*}
T=a+bP
\end{equation*}

At zero absolute pressure:
\begin{equation*}
T=bP
\end{equation*}

The ideal-gas temperature scale agrees with the thermodynamic scale over its usable range, but fails at extreme temperatures because of effects such as condensation, dissociation, and ionization.

\mysubsection{}{Temperature Conversion}

\begin{equation*}
\boxed{T(K)=T(^\circ\mathrm{C})+273.15}
\end{equation*}

\begin{equation*}
\boxed{T(R)=T(^\circ\mathrm{F})+459.67}
\end{equation*}

\begin{equation*}
\boxed{T(R)=1.8T(K)}
\end{equation*}

\begin{equation*}
\boxed{T(^\circ\mathrm{F})=1.8T(^\circ\mathrm{C})+32}
\end{equation*}

For engineering approximations, $273.15\approx273$ and $459.67\approx460$.

\mysubsection{}{Temperature Differences}

\begin{equation*}
\boxed{\Delta T(K)=\Delta T(^\circ\mathrm{C})}
\end{equation*}

\begin{equation*}
\boxed{\Delta T(R)=\Delta T(^\circ\mathrm{F})}
\end{equation*}

\begin{equation*}
\boxed{1~K=1^\circ\mathrm{C}=1.8~R=1.8^\circ\mathrm{F}}
\end{equation*}

\textbf{Critical distinction:}
\begin{itemize}
    \item Temperature \textbf{itself} in thermodynamic equations $\rightarrow$ use absolute scale.
    \item Temperature \textbf{difference} $\rightarrow$ K and $^\circ\mathrm{C}$ have identical magnitude.
\end{itemize}

%=========================================================
\mysection{1.9}{Pressure}{1.9}

\mysubsection{}{Definition and Units}

\begin{equation*}
\boxed{P=\frac{F_n}{A}}
\end{equation*}

Pressure is the normal force per unit area and is a scalar quantity.

\begin{equation*}
\boxed{1~\mathrm{Pa}=1~\mathrm{N/m^2}}
\end{equation*}

\begin{center}
\begin{tabular}{|c|c|}
\hline
$1~\mathrm{kPa}$ & $10^3~\mathrm{Pa}$ \\
\hline
$1~\mathrm{MPa}$ & $10^6~\mathrm{Pa}$ \\
\hline
$1~\mathrm{bar}$ & $10^5~\mathrm{Pa}=100~\mathrm{kPa}$ \\
\hline
$1~\mathrm{atm}$ & $101325~\mathrm{Pa}=101.325~\mathrm{kPa}$ \\
\hline
$1~\mathrm{kgf/cm^2}$ & $9.807\times10^4~\mathrm{Pa}$ \\
\hline
$1~\mathrm{psi}$ & $1~\mathrm{lbf/in^2}$ \\
\hline
\end{tabular}
\end{center}

Useful:
\begin{equation*}
1~\mathrm{atm}=1.01325~\mathrm{bar}=14.696~\mathrm{psi}
\end{equation*}

\begin{equation*}
1~\mathrm{kgf/cm^2}
=0.9807~\mathrm{bar}
=0.9679~\mathrm{atm}
=14.223~\mathrm{psi}
\end{equation*}

\mysubsection{}{Absolute, Gage, and Vacuum Pressure}

\begin{equation*}
\boxed{P_{\mathrm{gage}}=P_{\mathrm{abs}}-P_{\mathrm{atm}}}
\end{equation*}

\begin{equation*}
\boxed{P_{\mathrm{vac}}=P_{\mathrm{atm}}-P_{\mathrm{abs}}}
\end{equation*}

Therefore:
\begin{equation*}
\boxed{P_{\mathrm{abs}}=P_{\mathrm{atm}}+P_{\mathrm{gage}}}
\end{equation*}

\begin{equation*}
\boxed{P_{\mathrm{abs}}=P_{\mathrm{atm}}-P_{\mathrm{vac}}}
\end{equation*}

\begin{center}
\begin{tabular}{|c|c|c|}
\hline
\textbf{Pressure} & \textbf{Reference} & \textbf{Meaning} \\
\hline
Absolute & Vacuum & Actual pressure \\
\hline
Gage & Atmosphere & Pressure relative to atmosphere \\
\hline
Vacuum & Atmosphere & Amount below atmospheric pressure \\
\hline
\end{tabular}
\end{center}

\textbf{Unless specified otherwise, thermodynamic pressure means absolute pressure.}

\mysubsection{}{Hydrostatic Pressure}

For a fluid at rest:
\begin{itemize}
    \item Same horizontal level in the same continuous fluid $\Rightarrow$ same pressure.
    \item Pressure increases with depth.
    \item Pressure acts normal to a surface.
\end{itemize}

For constant density:
\begin{equation*}
\boxed{\Delta P=P_2-P_1=-\rho g\Delta z}
\end{equation*}

Equivalently:
\begin{equation*}
\boxed{P_{\mathrm{below}}
=P_{\mathrm{above}}+\rho g|\Delta z|}
\end{equation*}

or
\begin{equation*}
\boxed{P_{\mathrm{below}}
=P_{\mathrm{above}}+\gamma_s|\Delta z|}
\end{equation*}

For variable density:
\begin{equation*}
\boxed{\frac{dP}{dz}=-\rho g}
\end{equation*}

\begin{equation*}
\boxed{
P_2-P_1=-\int_1^2\rho g\,dz
}
\end{equation*}

For gases, pressure variation with elevation is usually negligible over small/moderate elevation changes.

\mysubsection{}{Pascal's Law}

Pressure applied to a confined fluid increases by the same amount throughout the fluid.

For connected pistons at the same level:
\begin{equation*}
P_1=P_2
\end{equation*}

\begin{equation*}
\boxed{\frac{F_1}{A_1}=\frac{F_2}{A_2}}
\end{equation*}

\begin{equation*}
\boxed{\frac{F_2}{F_1}=\frac{A_2}{A_1}}
\end{equation*}

\begin{equation*}
\boxed{\mathrm{IMA}=\frac{A_2}{A_1}}
\end{equation*}

Larger piston $\Rightarrow$ larger output force.

%=========================================================
\mysection{1.10}{Pressure Measurement Devices}{1.10}

\mysubsection{}{Barometer}

Measures atmospheric pressure.

\begin{equation*}
\boxed{P_{\mathrm{atm}}=\rho gh}
\end{equation*}

where $h$ is mercury-column height.

\begin{itemize}
    \item Column height is independent of tube length and cross-sectional area when capillary effects are negligible.
    \item Atmospheric pressure decreases with elevation.
\end{itemize}

\textbf{Standard atmosphere:}
\begin{equation*}
\boxed{1~\mathrm{atm}=760~\mathrm{mmHg}=760~\mathrm{torr}}
\end{equation*}

\begin{equation*}
1~\mathrm{torr}\approx133.3~\mathrm{Pa}
\end{equation*}

A water column of approximately $10.3~\mathrm{m}$ produces standard atmospheric pressure.

\mysubsection{}{Manometer}

A manometer determines pressure differences from liquid-column height.

Basic relation:
\begin{equation*}
\boxed{\Delta P=\rho gh}
\end{equation*}

\begin{center}
\begin{tabular}{|c|c|}
\hline
\textbf{Movement} & \textbf{Pressure change} \\
\hline
Downward through fluid & $+\rho gh$ \\
\hline
Upward through fluid & $-\rho gh$ \\
\hline
Horizontal in same continuous fluid & $0$ \\
\hline
\end{tabular}
\end{center}

For a manometer open to atmosphere:
\begin{equation*}
\boxed{P_2=P_{\mathrm{atm}}+\rho gh}
\end{equation*}

\textbf{Inclined manometer:} slanted tube increases liquid displacement for a given pressure difference $\Rightarrow$ higher resolution for small pressure differences.

\mysubsection{}{Multifluid Manometers}

For stacked immiscible fluids:
\begin{equation*}
\boxed{
P_{\mathrm{bottom}}
=
P_{\mathrm{top}}
+\rho_1gh_1+\rho_2gh_2+\rho_3gh_3
}
\end{equation*}

If all densities are equal:
\begin{equation*}
P_{\mathrm{bottom}}
=
P_{\mathrm{top}}
+\rho g(h_1+h_2+h_3)
\end{equation*}

\textbf{Analysis rules:}
\begin{enumerate}
    \item Downward $\rightarrow$ add $\rho gh$.
    \item Upward $\rightarrow$ subtract $\rho gh$.
    \item Same elevation in same continuous static fluid $\rightarrow$ equal pressure.
\end{enumerate}

\mysubsection{}{Differential Manometer}

Measures pressure difference between two specified points.

For the arrangement considered:
\begin{equation*}
\boxed{
P_1-P_2=(\rho_2-\rho_1)gh
}
\end{equation*}

where $\rho_1$ = flowing-fluid density, $\rho_2$ = manometer-fluid density.

For gases, $\rho_1\ll\rho_2$:
\begin{equation*}
\boxed{P_1-P_2\approx\rho_2gh}
\end{equation*}

%=========================================================
\mysection{1.11}{Problem-Solving Technique}{1.11}

\mysubsection{}{Seven-Step Procedure}

\begin{center}
\begin{tabular}{|c|l|}
\hline
\textbf{Step} & \textbf{Action} \\
\hline
1 & \textbf{Problem statement}: identify given data, unknowns, objective. \\
\hline
2 & \textbf{Schematic}: draw system; show dimensions, properties, interactions, unknowns. \\
\hline
3 & \textbf{Assumptions}: state and justify approximations. \\
\hline
4 & \textbf{Physical laws}: identify system/control volume and apply relevant laws. \\
\hline
5 & \textbf{Properties}: obtain required properties from relations/tables. \\
\hline
6 & \textbf{Calculations}: substitute, solve, check units and significant digits. \\
\hline
7 & \textbf{Verification}: check physical reasonableness, assumptions, validity, limitations. \\
\hline
\end{tabular}
\end{center}

\mysubsection{}{Exam/Problem-Solving Checklist}

\begin{itemize}
    \item Define the system/control volume before applying a law.
    \item State assumptions explicitly.
    \item Keep units throughout calculations.
    \item Check dimensional consistency.
    \item Retain sufficient intermediate digits.
    \item Check whether the final result is physically reasonable.
    \item Do not apply results outside the assumptions under which they were obtained.
\end{itemize}

\textbf{Engineering solutions:} neatness, organization, completeness, and clarity help reveal errors. Not every simple problem requires every step to be written explicitly.

%=========================================================
\newpage
\begin{center}
{\headingfont\Huge Significant Digits}
\par
\vspace{4pt}
\hrule
\end{center}

\mysubsection{}{Core Rules}

\begin{itemize}
    \item Significant digits represent meaningful precision in measured data.
    \item Final answer cannot be more precise than the least precise input.
    \item Keep full calculator precision during intermediate calculations.
    \item Round only the final answer.
    \item More digits do \textbf{not} imply greater accuracy.
    \item Experimental uncertainty limits meaningful precision.
\end{itemize}

Example:
\begin{equation*}
V=3.75~\mathrm{L},
\qquad
\rho=0.845~\mathrm{kg/L}
\end{equation*}

\begin{equation*}
m=V\rho
=3.75(0.845)
=3.16875~\mathrm{kg}
\end{equation*}

Since both inputs have three significant digits:
\begin{equation*}
\boxed{m=3.17~\mathrm{kg}}
\end{equation*}

\mysubsection{}{Rounding Uncertainty}

For a value rounded to a given decimal place, maximum rounding uncertainty is half the value of the least significant place.

For $3.75$:
\begin{equation*}
\text{uncertainty}
=
\pm\frac{0.01}{2}
=
\pm0.005
\end{equation*}

Therefore:
\begin{equation*}
\boxed{
3.745\le x<3.755
}
\end{equation*}

General rule:
\begin{equation*}
\boxed{
\text{rounding uncertainty}
=
\pm\frac{\text{value of least significant digit}}{2}
}
\end{equation*}

\mysubsection{}{Engineering Approximations}

Small intentional approximations may be acceptable when their error is negligible relative to other uncertainties.

Example:
\begin{equation*}
\rho_{\mathrm{water},75^\circ\mathrm{C}}
\approx975~\mathrm{kg/m^3}
\rightarrow
1000~\mathrm{kg/m^3}
\end{equation*}

Approximate error:
\begin{equation*}
\frac{1000-975}{975}\times100\%\approx2.5\%
\end{equation*}

Thus, reasonable rounding and approximation are acceptable when justified by the required accuracy.

\end{document}